\documentclass[11pt,a4paper]{article}

\usepackage[margin=2.2cm]{geometry}
\usepackage{booktabs}
\usepackage{longtable}
\usepackage{tabularx}
\usepackage{array}
\usepackage{hyperref}
\usepackage{enumitem}
\usepackage{listings}
\usepackage{xcolor}
\usepackage{graphicx}
\usepackage{amsmath}

\hypersetup{
  colorlinks=true,
  linkcolor=black,
  urlcolor=blue
}

\lstset{
  basicstyle=\ttfamily\small,
  breaklines=true,
  frame=single,
  backgroundcolor=\color{gray!8},
  columns=fullflexible
}

\title{AgriGuard Hybrid Anomaly Detection Model\\
\large Live GPS Simulation + Rules + Isolation Forest}
\author{AgriGuard Project Notes}
\date{15 August 2026}

\begin{document}
\maketitle

\begin{abstract}
AgriGuard's livestock-theft detector is a \textbf{hybrid} system: hard rules fire first for
clear events (geofence breach, high speed, night movement out of zone), then an
\textbf{Isolation Forest} (with a heading assist) catches subtler
\textbf{Erratic Movement} such as zig-zag paths that still look normal on speed and fence alone.
Live tracking feeds the detector one simulated GPS collar ping per active animal about every
5~seconds. There are no physical collars; the backend invents constrained movement to mimic
a star (hub-and-spoke) reporting topology.
\end{abstract}

\tableofcontents
\vspace{1em}

%----------------------------------------------------------------------
\section{Big picture}
%----------------------------------------------------------------------

\begin{lstlisting}
Farmer browser (tracking.html)
        |  every 5s: POST /api/tracking/live-tick
        v
Flask invents next collar ping (simulation)
        |
        v
Rules in live-tick  -->  Geofence Breach / High Speed / Night Movement
        | (if none of those)
        v
AnomalyDetector.predict()  -->  Isolation Forest + heading assist
        |
        v
gps_tracking + alerts (+ email for breach/speed only)
\end{lstlisting}

\noindent
Each animal is scored \textbf{independently} (per animal, per ping)---not by a herd vote.

%----------------------------------------------------------------------
\section{What one ping looks like}
%----------------------------------------------------------------------

Each live-tick, for every \textbf{Active} animal of the logged-in farmer:

\begin{enumerate}[leftmargin=*]
  \item Start from last latitude/longitude (or place inside fence/zone).
  \item Choose behaviour: rest / graze / walk / rare anomaly.
  \item Compute new latitude, longitude, and speed.
  \item Update heading and the rolling \textbf{mean turn size} over recent ticks.
  \item Classify the ping (rules, then ML).
  \item Insert into \texttt{gps\_tracking}; update \texttt{animals.last\_latitude/longitude}.
  \item If anomalous, create an \texttt{alerts} row.
\end{enumerate}

%----------------------------------------------------------------------
\section{Simulation behaviour}
%----------------------------------------------------------------------

The UI calls \texttt{POST /api/tracking/live-tick} about every \textbf{5 seconds}
while Live Tracking is on.

\begin{center}
\begin{tabular}{@{}llp{7.2cm}@{}}
\toprule
\textbf{Behaviour} & \textbf{Chance} & \textbf{What it does} \\
\midrule
Rest    & $\sim$25\% & Almost no move, $\approx 0$~km/h \\
Graze   & $\sim$55\% & Small step (cattle $\sim$3--15~m), slow \\
Walk    & $\sim$15\% & Larger step, still under $\sim$4--4.5~km/h \\
Anomaly & $\sim$5\%  & High Speed, Geofence Breach, Erratic,
                       or Night (if Night Mode + assigned zone) \\
\bottomrule
\end{tabular}
\end{center}

\paragraph{Normal moves.}
Forced to stay \textbf{inside} the farm geofence (polygon preferred; circle fallback).

\paragraph{Erratic Movement (simulation).}
\begin{itemize}[leftmargin=*]
  \item Starts a multi-tick zig-zag streak (2--5 further ticks).
  \item Each tick: large left/right turn ($\sim$90$^\circ$--160$^\circ$).
  \item Speed stays \textbf{under 8~km/h} and path stays \textbf{inside the fence}.
  \item Therefore rules P1--P3 do not fire; ML / heading must catch it.
\end{itemize}

Heading memory keeps successive steps roughly continuous (not teleporting).

%----------------------------------------------------------------------
\section{The four ML features}
%----------------------------------------------------------------------

Every Isolation Forest call uses a 4-dimensional feature vector:

\begin{center}
\begin{tabular}{@{}clp{8cm}@{}}
\toprule
\# & \textbf{Feature} & \textbf{Meaning} \\
\midrule
1 & \texttt{speed} & Movement speed in km/h \\
2 & \texttt{hour} & Hour of day 0--23
  (simulation: Night Mode $\rightarrow$ 22; day $\rightarrow$ 12) \\
3 & \texttt{distance\_from\_center} & Metres from farm centre
  (helper for circle/ML geometry) \\
4 & \texttt{heading\_variance} & Rolling \textbf{mean} of recent turn sizes
  (radians)---despite the name \\
\bottomrule
\end{tabular}
\end{center}

\subsection{Heading feature (mean turn size)}
\begin{itemize}[leftmargin=*]
  \item Each tick: absolute turn angle between previous and new heading.
  \item Keep the last \textbf{5} turns per animal.
  \item Feature value $=$ \textbf{mean} of those turns (not standard deviation).
\end{itemize}

\paragraph{Why mean, not standard deviation?}
\begin{itemize}[leftmargin=*]
  \item Grazing drifts a little $\Rightarrow$ mean $\approx$ \textbf{0.1--0.4}.
  \item Erratic keeps making huge turns $\Rightarrow$ mean $\approx$ \textbf{1.8--2.8}.
  \item Std-dev of ``all large but similar turns'' can stay low and hide zig-zag.
\end{itemize}

(The parameter is still named \texttt{heading\_variance} in code for compatibility;
the numeric value is mean turn size.)

%----------------------------------------------------------------------
\section{Detection priority (hybrid stack)}
%----------------------------------------------------------------------

The same priority appears in live-tick classification and in
\texttt{AnomalyDetector.predict()}:

\subsection{P1 --- Geofence Breach}
\begin{itemize}[leftmargin=*]
  \item Live path: point outside farm polygon/circle (\texttt{is\_inside\_farm}).
  \item Detector fallback: distance $>$ radius $\times 1.1$.
\end{itemize}
Most objective signal: the animal left the farm.

\subsection{P2 --- High Speed}
Speed $> 8$~km/h. Physics-based: faster than normal graze/walk.

\subsection{P3 --- Night Movement}
Night \textbf{and} outside the animal's assigned \textbf{zone}
(not merely ``any night movement''). In the demo, Night Mode UI drives night hour.

\subsection{P4 --- ML + heading assist}
Only if P1--P3 did not fire:
\begin{enumerate}[leftmargin=*]
  \item Run Isolation Forest on the four features.
  \item If forest predicts outlier (\texttt{predict == -1}) $\rightarrow$
        \textbf{Erratic Movement}.
  \item Else if mean turn $\geq 1.2$ $\rightarrow$ \textbf{Erratic Movement}
        (demo-reliable assist).
\end{enumerate}

So Erratic is the main ML job: a weird path that still looks fine on speed and fence.

%----------------------------------------------------------------------
\section{Isolation Forest (the ML component)}
%----------------------------------------------------------------------

\begin{center}
\begin{tabular}{@{}ll@{}}
\toprule
\textbf{Property} & \textbf{Value} \\
\midrule
Library & \texttt{sklearn.ensemble.IsolationForest} \\
Trees (\texttt{n\_estimators}) & 100 \\
Contamination & 0.15 ($\sim$15\% of space treated as outliers) \\
Random state & 42 (reproducible) \\
Saved model & \texttt{Backend/ml/models/anomaly\_model\_v2.pkl} \\
Source module & \texttt{Backend/ml/anomaly\_detector.py} \\
\bottomrule
\end{tabular}
\end{center}

\subsection{Idea}
Isolation trees randomly partition the feature space. Outliers are easier to isolate
and receive more anomalous scores. Fitting is \textbf{unsupervised} (no class labels).

\subsection{Synthetic training set (1000 samples)}
\begin{itemize}[leftmargin=*]
  \item \textbf{700 normal:} daytime, slow speed, inside fence, small mean turns.
  \item \textbf{300 anomaly-shaped} rows (to enrich the space):
        high speed / night hours / large distance / \textbf{high mean turn (erratic)}.
\end{itemize}

\subsection{Predict outputs}
\begin{itemize}[leftmargin=*]
  \item $+1$ --- looks normal.
  \item $-1$ --- outlier $\rightarrow$ labelled \textbf{Erratic Movement} in the API.
  \item \texttt{score\_samples} --- anomaly score (more negative $\approx$ more anomalous).
\end{itemize}

\subsection{Caveat}
Putting anomaly-like rows into Isolation Forest training is somewhat unusual
(the algorithm assumes mostly normal data). In practice, rules catch breach / speed / night;
the forest plus assist mainly separate \textbf{smooth vs zig-zag} via the heading feature.

\subsection{Startup}
\texttt{ensure\_ready()} \textbf{loads} \texttt{anomaly\_model\_v2.pkl} if present;
it trains only when the file is missing (does not retrain on every Flask boot).

%----------------------------------------------------------------------
\section{End-to-end example}
%----------------------------------------------------------------------

Animal zig-zagging inside the fence at about 4~km/h:

\begin{enumerate}[leftmargin=*]
  \item Simulation picks or continues an Erratic streak $\rightarrow$ large turn, short step.
  \item Mean turn over recent ticks climbs toward $\sim$2.
  \item Live rules: inside fence, speed $< 8$, not night/out-of-zone $\rightarrow$ pass.
  \item Call
        \texttt{predict(speed=4, hour=12, distance=300, heading\_variance=2.1)}.
  \item Detector rules P1--P3 also pass.
  \item Forest and/or assist ($\geq 1.2$) $\rightarrow$ \textbf{Erratic Movement}.
  \item Insert \texttt{gps\_tracking} with \texttt{is\_anomaly=true}.
  \item Create alert; \textbf{no email} (email only for Breach / High Speed).
\end{enumerate}

\paragraph{UI note.}
Map marker red $=$ latest ping's \texttt{is\_anomaly} only (clears on next normal ping).
An alert stays Active until the farmer resolves it.

%----------------------------------------------------------------------
\section{Alert types}
%----------------------------------------------------------------------

\begin{center}
\begin{tabular}{@{}lp{9cm}@{}}
\toprule
\textbf{Type} & \textbf{Typical cause} \\
\midrule
Geofence Breach   & Outside farm fence \\
High Speed        & Speed $> 8$~km/h \\
Night Movement    & Night Mode on and outside assigned zone \\
Erratic Movement  & Zig-zag / high mean turn (ML + assist) \\
\bottomrule
\end{tabular}
\end{center}

%----------------------------------------------------------------------
\section{Networking topology}
%----------------------------------------------------------------------

AgriGuard mimics a \textbf{star (hub-and-spoke)} topology:

\begin{itemize}[leftmargin=*]
  \item \textbf{Hub:} AgriGuard Flask API (stores GPS, runs rules + ML).
  \item \textbf{Spokes:} each animal's simulated collar reports independently.
  \item Not mesh and not peer-to-peer: animals do not communicate with each other.
  \item Transport in the demo: REST polling (\texttt{live-tick}), not WebSockets.
\end{itemize}

%----------------------------------------------------------------------
\section{Key files}
%----------------------------------------------------------------------

\begin{center}
\begin{tabular}{@{}lp{8.5cm}@{}}
\toprule
\textbf{Path} & \textbf{Role} \\
\midrule
\texttt{Backend/ml/anomaly\_detector.py}
  & Rules + Isolation Forest v2 \\
\texttt{Backend/ml/models/anomaly\_model\_v2.pkl}
  & Saved forest \\
\texttt{Backend/app.py} (live-tick section)
  & Simulation, heading mean, classify, DB, alerts \\
\texttt{Frontend/farmer/tracking.html}
  & Map UI, Live Tracking, Night Mode \\
\bottomrule
\end{tabular}
\end{center}

%----------------------------------------------------------------------
\section{One-line summary}
%----------------------------------------------------------------------

\textbf{Simulate realistic GPS $\rightarrow$ rules for clear theft signals $\rightarrow$
Isolation Forest + mean heading for erratic paths that look ``normal'' on speed and fence alone.}

\end{document}
